\renewcommand{\arraystretch}{1.3}
\captionsetup[table]{justification=raggedright,singlelinecheck=false}
\begin{table}[H]
\caption{Distribusi karakteristik \textit{node} berdasarkan kedalaman penelusuran graf}
\label{tbl:hop_justification}
\centering
\begin{tabular}{| c | p{0.20\textwidth} | p{0.30\textwidth} | p{0.25\textwidth} |}
\hline
\textbf{Kedalaman} & \textbf{Karakteristik \textit{Node}} & \textbf{Contoh} & \textbf{Penggunaan} \\
\hline
1--2 & Spesifik \textit{resource} (dibagikan 2--3 \textit{resource}) & \textit{DeploymentSpec}, \textit{ServiceSpec} & Tipe \textit{explain}, \textit{followup} \\
\hline
3 & \textit{Field} tingkat kontainer & \textit{image}, \textit{ports}, \textit{env}, \textit{resources} & Tipe \textit{generate\_yaml}, \textit{trace\_relationship} \\
\hline
4+ & Tipe generik (dibagikan 19--136 \textit{resource}) & \textit{Quantity}, \textit{IntOrString}, \textit{LocalObjectReference} & \textbf{Tidak digunakan} (menurunkan presisi) \\
\hline
\end{tabular}
\end{table}
